\section*{Physique (13 points)}

\section*{Exercice 1 (3,5 points)~: Transformations nucléaires}
L'or $\ce{_{79}^{198}Au}$ est un radioisotope de l'or, il se désintègre en
mercure $\ce{_{Z}^{A}Hg}$ par émission $\beta^{-}$. Cette radioactivité est
notamment utilisée dans le traitement contre certains cancers ou pour traiter
d'autres maladies.

Pour évaluer l'efficacité thérapeutique de l'or 198, une équipe de recherche a
porté son attention sur un groupe de souris souffrantes d'un cancer de la
prostate sévère en leurs injectant une dose d'or $\ce{_{79}^{198}Au}$.

\begin{donnees}
    \begin{table}[H]
        \centering
        \setlength{\arrayrulewidth}{0.8pt}
        \renewcommand{\arraystretch}{1.5}
        \begin{tabularx}{\textwidth}{|p{4cm}|C|C|C|C|C|}
            \hline
            Noyau ou particule & $\ce{_{79}^{198}Au}$ & $\ce{_{Z}^{A}Hg}$ &
            Proton & Neutron & Électron \\
            \hline
            Masse en (\unit{\u}) & $197,968225$ & $197,966769$ & $1,007276$ &
            $1,008665$ & $\num{5.49e-4}$ \\
            \hline
            \multicolumn{4}{|l|}{\raisebox{2mm}[9mm]{
                Énergie de liaison par nucléon~:
                $\frac{E_{\ell}}{A}\left(\ce{_{Z}^{A}Hg}\right)=
                \qty{8.4}{\MeV\per\mathrm{nucléon}}$}} &
            \multicolumn{2}{c|}{\raisebox{2mm}[9mm]{
                $\qty{1}{\u}=\qty{931.5}{\U}$}} \\
            \hline
        \end{tabularx}
    \end{table}
\end{donnees}

\begin{enumerate}
    \item Écrire l'équation de désintégration de l'or $\ce{_{79}^{198}Au}$ en
          précisant les valeurs de $A$ et $Z$.
    \item Vérifier que la valeur de l'énergie de liaison
          $E_{\ell}(\ce{_{79}^{198}Au})$ est égale à $\qty{1525.53}{\MeV}$.
    \item Justifier que le noyau $\ce{_{Z}^{A}Hg}$ est plus stable que la noyau
          $\ce{_{79}^{198}Au}$.
    \item Calculer, en unité $\unit{\MeV}$, l'énergie libérée
          $E_{\text{libérée}}=|\Delta E|$ par la réaction de désintégration d'un
          noyau d'or 198.
\end{enumerate}



\section*{Exercice 2 ( 5.5 points)~: Dipôle RC - Circuit LC}
La charge et la décharge d'un condensateur sont deux phénomènes électriques dont
chacun est régi par une équation décrivant l'évolution de la tension aux bornes
du condensateur. Lorsque le condensateur est associé à une bobine, le circuit LC
obtenu peut être siège d'oscillations électriques selon un régime donné.

Cet exercice vise~:
\begin{itemize}
    \item l'étude de la réponse d'un dipôle RC à un échelon de tension
          ascendant~;
    \item l'étude énergétique d'un circuit oscillant LC.
\end{itemize}

\subsection*{Partie 1~: Réponse d'un dipôle RC à un échelon de tension ascendant}
Un générateur de tension, de force électromotrice $E$ alimente un conducteur
ohmique de résistance $R=\qty{100}{\ohm}$ et un condensateur de capacité $C$.

Un oscilloscope numérique est utilisé pour suivre l'évolution temporelle des
deux tensions en voie~1 et en voie~2 (Fig.~\ref{fig:1}).

Le condensateur étant préalablement déchargé. À l'instant ${t_{0}=0}$, on ferme
l'interrupteur $\mathrm{K}$. Les évolutions temporelles des deux tensions sont
traduites par les courbes données par la Fig.~\ref{fig:2}.
\begin{minipage}{0.48\linewidth}
    \begin{figure}[H]
        \centering
        \includegraphics[width=0.7\textwidth]{1}
        \caption{}
        \label{fig:1}
    \end{figure}
\end{minipage}
\hfill
\begin{minipage}{0.48\linewidth}
    \begin{figure}[H]
        \centering
        \includegraphics[width=0.9\textwidth]{2}
        \caption{}
        \label{fig:2}
    \end{figure}
\end{minipage}

\begin{enumerate}
    \item Recopier sur votre copie le schéma de la Fig.~\ref{fig:1} et
          représenter, en convention récepteur, la tension $u_{C}$ aux bornes du
          condensateur et la tension $u_{R}$ aux bornes du conducteur ohmique.
    \item Parmi les courbes \ding{182} et \ding{183} de la Fig.~\ref{fig:2},
          quelle est celle qui correspond à la tension $u_{C}$~? Justifier.
    \item Montrer que l'équation différentielle vérifiée par la tension $u_{C}$
          s'écrit $R\cdot C\cdot\frac{\mathrm{d}u_{C}}{\mathrm{d}t}+u_{C}=E$.
    \item En admettant que $u_{c}=E\cdot\bigl(1-e^{-\frac{t}{\tau}}\bigr)$ est
          solution de l'équation différentielle, déterminer l'expression de
          $\tau$ en fonction des paramètres du circuit.
    \item Calculer la valeur de $C$.
    \item Recopier, sur votre copie, le numéro de la question et choisir la
          lettre correspondante à la proposition vraie. L'expression de
          l'intensité $i(t)$ du courant qui traverse le circuit est~:

          \begin{minipage}{\linewidth}
          \begin{table}[H]
              \centering
              \setlength{\arrayrulewidth}{0.8pt}
              \renewcommand{\arraystretch}{1.5}
              \begin{tabularx}{\textwidth}{|
                  >{\arraybackslash\centering}p{6mm}|C|
                  >{\arraybackslash\centering}p{6mm}|C|}
              \hline
              $\mathrm{A}$ & $i(t)=0,06\cdot e^{-3333,3\cdot t}~(\unit{\A})$
              &
              $\mathrm{B}$ & $i(t)=-0,06\cdot e^{-3333,3\cdot t}~(\unit{\A})$ \\
              \hline
              $\mathrm{C}$ & $i(t)=0,6\cdot e^{-3333,3\cdot t}~(\unit{\A})$
              &
              $\mathrm{D}$ & $i(t)=6\cdot e^{-3333,3\cdot t}~(\unit{\A})$ \\
              \hline
              \end{tabularx}
          \end{table}
          \end{minipage}
\end{enumerate}




\section*{Partie 2~: Étude énergétique d'un circuit oscillant LC}
On considère le circuit de la
Fig.~\ref{fig:2bac_svt_national_2024_rattrapage_phys_ex2_3} tel que~:
$E=\qty{10}{\V}$ et $C=\qty{3}{\uF}$.

La résistance du générateur et celle de la bobine sont supposées négligeables.

On place l'interrupteur $\mathrm{K}$ en position~(1). Lorsque le condensateur
est totalement chargé, on bascule $\mathrm{K}$ en position~(2) à l'instant
$t_{0}=0$.

La Fig.~\ref{fig:2bac_svt_national_2024_rattrapage_phys_ex2_4} donne les
variations de l'une des formes d'énergie~: l'énergie électrique
$\mathcal{E}_{e}$ emmagasinée dans le condensateur, l'énergie magnétique
$\mathcal{E}_{m}$ emmagasinée dans la bobine ou l'énergie totale $\mathcal{E}$
du circuit.

\begin{minipage}{0.48\linewidth}
    \begin{figure}[H]
        \centering
        \begin{circuitikz}
            \newdimen\d \d=3cm
            \draw (0,0) node[cute spdt mid,rotate=90] (Sw) {};
            \node[above] at (Sw.out 1) {(1)};
            \node[above] at (Sw.out 2) {(2)};
            \node[above=9mm] at (Sw.in) {$\mathrm{K}$};
            \draw (Sw.out 1)
                 to[short] ++(-.8\d,0)
                 to[vsource,v_<=$E$,voltage shift=0.5] ++(0,-\d)
                 to[short] ++(2\d,0) coordinate(M)
                 to[L,l_={$L$},mirror,voltage shift=3.5,name=C] ++(0,\d)
                 to[short] (Sw.out 2);
            \draw (Sw.in) to[C,l_=$C$,i>^=$i$] (Sw.in |- M);
        \end{circuitikz}
    \caption{}
    \label{fig:2bac_svt_national_2024_rattrapage_phys_ex2_3}
    \end{figure}
\end{minipage}
\hfill
\begin{minipage}{0.48\linewidth}
    \begin{figure}[H]
        \centering
        \includegraphics[width=.7\textwidth]{2bac_svt_national_2024_rattrapage_phys_ex2_4}
        \caption{}
        \label{fig:2bac_svt_national_2024_rattrapage_phys_ex2_4}
    \end{figure}
\end{minipage}

\begin{enumerate}
    \item Donner, en justifiant, la forme d'énergie correspondante à la courbe
          représentée dans la
          Fig.~\ref{fig:2bac_svt_national_2024_rattrapage_phys_ex2_4}.
    \item Expliquer de point de vue énergétique le régime d'oscillations
          électriques obtenu.
    \item Déterminer graphiquement la valeur de la période propre $T_{0}$ des
          oscillations.
    \item Déterminer la valeur de l'inductance $L$ de la bobine.
    \item Déterminer la valeur de l'énergie magnétique $\mathcal{E}_{m}$
          emmagasinée dans la bobine à l'instant $t=\frac{3 T_{0}}{4}$.
\end{enumerate}



\section*{Exercice 3 ( 4 points)~: Mouvement des systèmes mécaniques modélisés}
Dans la vie courante, un certain nombre de systèmes mécaniques sont utilisés et
modélisés. Les mouvements de ces systèmes peuvent être identifís à la suite
d'une étude dynamique, ce qui permet de déterminer des grandeurs (dynamiques et
cinématiques) caractérisant ces mouvements.

On se propose dans cet exercice d'étudier deux exemples de systèmes mécaniques
modélisés.

\subsection{Étude du mouvement d'un skeletoneur}
Le skeleton est un sport d'hiver qui se pratique individuellement sur une
planche où le skeletoneur se place à plat ventre, la tête devant. L'objectif est
de parcourir la piste le plus rapidement possible. On se propose d'étudier le
mouvement du centre d'inertie $G$ d'un skeletoneur, de masse $m$, sur une partie
d'une piste rectiligne inclinée d'un angle $\alpha$ par rapport à l'horizontal.
L'étude est réalisée dans un repère ($O,\ex,\ey$ ) liée à la Terre
supposé galiléen.

\begin{minipage}{0.65\linewidth}
    On choisit l'instant de passage de $G$ en $O$ comme origine de temps
    ($t_{0}=0$). On repère la position de $G$ à un instant $t$ par son abscisse
    $x_{G}$ dans ce repère (Figure~\ref{fig:skeletoneur}).
    
    Le skeletoneur arrive au point de départ $O$ avec la vitesse
    $\vec{v}_{0}=v_{0}\ex$, il passe ensuite par la position $A$ avec une
    vitesse $\vec{v}_{A}$.
    
    Le contact avec le plan incliné se fait avec frottement modélisé par une
    force $\vec{f}$ parallèle au plan, de sens opposée au sens de mouvement et
    d'intensité constante $f$.
    
    \begin{donnees}
        \begin{itemize}
            \item $g=\qty{9,8}{\a}$~; $\alpha=\ang{10}$~; $m=\qty{80}{kg}$~;
                  $v_{A}=\qty{30}{\v}$.
        \end{itemize}
    \end{donnees}
\end{minipage}
\hfill
\begin{minipage}{0.34\linewidth}
    \begin{figure}[H]
        \centering
        \includegraphics[width=\textwidth]{2025-05-17_014726-}
        \caption{}
        \label{fig:skeletoneur}
    \end{figure}
\end{minipage}

\begin{minipage}{0.65\linewidth}
\begin{enumerate}
    \item En appliquant la deuxième loi de Newton, montrer que l'accélération de
          $G$ s'écrit~: $a_{G}=g \sin \alpha-\frac{f}{m}$.
    \item La courbe de la figure~(\ref{fig:v_de_t}) donne l'évolution de la
          vitesse $v$ de $G$ au cours du temps.
          \begin{enumerate}
              \item Déterminer la valeur de $a_{G}$.
              \item Déduire, en justifiant, la nature du mouvement de $G$.
          \end{enumerate}   
    \item Calculer la valeur de $f$.
    \item Déterminer la valeur de l'instant $t_{A}$ de passage de $G$ par $A$.
    \item Calculer la distance $OA$.
\end{enumerate}
\end{minipage}
\hfill
\begin{minipage}{0.34\linewidth}
    \begin{figure}[H]
        \centering
        \includegraphics[width=\linewidth]{2025-05-17_015326-}
        \caption{}
        \label{fig:v_de_t}
    \end{figure}
\end{minipage}



\section*{Partie 2~: Étude du mouvement du système oscillant (Solide - Ressort)}
\begin{minipage}{0.65\linewidth}
    On considère le système oscillant (Solide-Ressort) où $m$ désigne la masse
    du solide et $K$ la raideur du ressort.

    On étudie le mouvement de $G$ dans un repère $(O,\ex)$ lié à la Terre
    supposé galiléen (Fig.~\ref{fig:solide_ressort}).

    On écarte $\mathrm{(S)}$ de sa position d'équilibre d'une distance $X_m$
    dans le
    sens opposé de $\ex$ puis on l'abandonne sans vitesse initiale à l'instant
    ($t_0=0$). Le solide $\mathrm{(S)}$ oscille sans frottement sur le plan
    horizontal.
\end{minipage}
\hfill
\begin{minipage}{0.34\linewidth}
    \begin{figure}[H]
        \centering
        \begin{tikzpicture}
    \node[inner sep=0pt,minimum width=4.5cm,minimum height=3mm,
        top color=black!70] (sol) {};
    % Axe xx'
    \coordinate (O) at ($(sol.south west)!.8!(sol.south east)$);
    \draw[-Latex] (sol.south west)
        node[below right,inner xsep=0pt] {$x^{\prime}$} -- (O)
        node[below=0.4mm] {$O$} -- (sol.south east) -- ++(0.7,0)
        node[below left,inner xsep=0pt,inner ysep=2mm] {$x$};
    % Fixed
    \node[inner sep=0pt,draw,minimum width=3mm,minimum height=7mm,
        anchor=south,pattern=north west lines,
        thick] (fix) at ([xshift=3mm,yshift=0.4pt]sol.north west) {};
    % Solide (S)
    \node[inner sep=0pt,draw,minimum width=7mm,minimum height=7mm,
        anchor=south,rounded corners=2pt,thick,label=$\mathrm{(S)}$,
        fill=lightgray] (S) at
        ([yshift=0.4pt]O |- sol.north) {};
        %([xshift=-6mm,yshift=0.4pt]sol.north east) {};
    \node[inner sep=1pt,circle,fill,
        label={[right=-0.8mm,font=\scriptsize]:$G$}] (G) at (S.center) {};
    \draw[{Rectangle[length=1pt,width=4pt]}-{Stealth[length=5pt]},
        thick] (O) -- ++(0.8,0)
        node[below left,inner xsep=0pt] {$\ex$};
    % Ressort
    \draw[decoration={
          coil,
          aspect=0.3,
          segment length=2.0mm,
          amplitude=2.4mm,
          pre length=1mm,
          post length=1mm},
          decorate,thick,black] (fix.east) -- (S.west);
\end{tikzpicture}

        \caption{}
        \label{fig:solide_ressort}
    \end{figure}
\end{minipage}

L'équation horaire du mouvement de $G$ est~:
${x(t)=\num{5e-2}\cdot\cos(4\cdot t+\pi)}$ (\unit{\meter}).
\begin{donnees}
    \begin{itemize}
        \item $m=\qty{1.25}{\kg}$
    \end{itemize}
\end{donnees}
\begin{enumerate}
    \item Déterminer la valeur de la période propre $T_0$.
    \item Déduire la valeur de la raideur $K$.
    \item Déterminer la valeur de $\dot{x}_G$ vitesse de $G$ au passage par la
          position d'équilibre pour la première fois dans le sens positif.
\end{enumerate}


